\section{Discussion}
\label{sec:discussion}

\subsection{Reframing the Alignment Tax on Diversity}
\label{sec:reframing}

Our results challenge the common narrative that alignment diminishes stylistic
capabilities. The base model has diverse latent distributions---this is
well-established by the pretraining literature and confirmed by the smooth
interpolation curves in our experiments. However, having a distribution and
being able to \emph{access} it on demand are fundamentally different
capabilities. The base model scores $3.74$ on style fidelity not because it
lacks stylistic knowledge, but because it lacks the instruction-following
mechanism needed to direct that knowledge toward a specific target.

This reframing aligns with the Superficial Alignment Hypothesis~\cite{zhou2023lima}
and the elasticity findings of \citet{ji2024resistance}. Alignment does not
rewrite the model's distributional structure; it adds a thin control layer.
The ``alignment tax'' on diversity documented by \citet{kirk2023rlhf} is real at
the level of \emph{unsolicited} generation---aligned models default to a narrower
range of outputs when not specifically instructed to diversify. But when
\emph{instructed} to produce a specific style, the aligned model draws on the same
distributional knowledge as the base model, and does so far more effectively.

\subsection{Distribution Arithmetic as a Diagnostic Tool}
\label{sec:diagnostic}

The smooth, monotonic interpolation curve (\figref{fig:interpolation}) serves as
evidence that alignment operates continuously on the base distribution. There is
no phase transition or discontinuity between base and aligned behavior---the two
are connected by a navigable path in logit space. This is consistent with the
log-linear structure exploited by \citet{zhou2024emulated} and the
continuous $\lambda$-interpolation of \citet{zhu2025adding}.

The concave shape of the curve is informative: most style improvement occurs
between $\interpcoeff = 0$ and $\interpcoeff = 0.5$, suggesting that the aligned
model's primary contribution is the instruction-following ``address'' rather than
new stylistic knowledge. The aligned model does not know more about Hemingway than
the base model does; it simply knows how to use what it knows when asked.

\subsection{Why Some Styles Are Harder}
\label{sec:style_difficulty}

The per-style analysis (\tabref{tab:perstyle}) reveals a pattern: ornate,
additive styles (Poe, Shakespeare, noir) respond most to alignment, while the
minimalist Hemingway style shows the smallest improvement. We hypothesize that
this reflects an asymmetry in the generation process. Ornate styles require
\emph{adding} distinctive features---gothic imagery, archaic vocabulary,
hard-boiled metaphors---which the instruction-following mechanism can reliably
trigger. Hemingway's style, by contrast, requires \emph{restraint}: shorter
sentences, simpler vocabulary, the omission of elaboration. ``Doing less'' is a
harder optimization target for autoregressive models, which are trained to predict
the next token, not to suppress tokens.

\subsection{Limitations}
\label{sec:limitations}

\para{Single model pair.}
We tested only \qwenbase and \qweninst. Results may differ for larger models,
different model families, or models aligned with different methods (PPO vs.\ DPO
vs.\ RLHF). Larger models may show different patterns due to greater distributional
capacity.

\para{LLM-as-judge bias.}
Style evaluation via \gptjudge is subjective and may carry systematic
biases---for example, preferring fluent, well-structured text, which inherently
advantages the aligned model. Human evaluation would provide a valuable complement,
though it introduces its own biases regarding style recognition.

\para{Benign styles only.}
Our seven target styles are all benign literary and functional styles. The more
consequential case for the alignment tax hypothesis involves styles that alignment
actively suppresses---toxic content, impersonation of real living persons, or other
safety-relevant registers. Distribution arithmetic may show qualitatively different
behavior in those settings.

\para{Sample size.}
With 2--3 samples per condition--task pair, statistical power for per-style
comparisons is limited. The aggregate results are robust ($N = 231$ total samples),
but individual style estimates should be interpreted cautiously.

\para{Prompt format confound.}
The base model received completion-style prompts while the aligned model used chat
templates. This is a necessary design choice---each model performs best with its
native format---but it means the aligned model's advantage may partly reflect
superior prompt formatting rather than alignment \emph{per se}.

\para{Logit-space only.}
We tested only logit-space distribution arithmetic, not weight-space model
averaging~\cite{lin2023mitigating}. Weight-space interpolation may exhibit
different properties, particularly for layer-specific blending strategies.
